\section{Further extremal deductions}\label{sec:deductions}

The formal witness also supports stronger questions than the first
existential optimum. We now distinguish the minimum for a \emph{fixed}
opening, the least possible mass when the target is reached, and the
largest score among the fastest games. Separately, binary mass barriers
yield policy-independent stochastic bounds and a useful restriction of
the unresolved maximum-score search. Only the results explicitly marked
as Lean-checked below are included in the extended formal development.

\subsection{Joint time--mass optima for every opening}

\begin{theorem}[Every-opening joint optimum; Lean-checked]\label{thm:allopenings}
Let $A$ be any ordinary two-tile opening of standard $4\times4$ 2048, and put
\[
 m(A)=\begin{cases}
 32{,}781,&\text{both opening tiles are 4s},\\
 32{,}782,&\text{otherwise}.
 \end{cases}
\]
There is a legal play from $A$ of exactly $m(A)$ rounds ending with a
$131072$ tile and total mass $131102$. Every legal play from $A$ reaching
that tile has at least $m(A)$ rounds and endpoint mass at least $131102$.
\end{theorem}

The quantifiers are stronger than in Theorem~\ref{thm:leanexact}:
\emph{for every opening there exists a continuation}. There are 480
labeled openings: 120 with two 4s, 240 mixed, and 120 with two 2s. They
form 75 square-symmetry classes. A finite bridge search finds, for each
opening, at most three preparatory rounds joining a symmetry of one of
the first four states of the original certificate. The length of the
remaining suffix is accounted for exactly, rather than adding a blanket
three-round overhead. A separately proved symmetry transport, every
bridge, and exhaustive coverage of the ordinary-opening definition are
checked in Lean. The new theorem is
\begin{quote}\small\ttfamily
Game2048.standard2048\_all\_openings\_exact.
\end{quote}
The constructive half has no axiom dependencies; the combined theorem
uses only \texttt{propext} and \texttt{Quot.sound}.

\begin{proof}[Mass lower bound and construction]
Continue to write $H_k(B)$ for the mass in tiles strictly greater than
$k$. A round $B\to B'$ with $k\ge2$ satisfies
\[
 H_{2k}(B')\le H_k(B),\qquad \mass(B')\ge\mass(B)+2.
\]
Thus the mass outside the increasing threshold grows by at least two
per round. Induction over any final $r$ rounds gives
\[
 H_{2^{r+1}}(B')+2r\le\mass(B').
\]
With $r=15$, a $131072$ tile contributes to $H_{65536}$, so the endpoint
mass is at least $131072+30=131102$. The existing move lower bound
ensures that these fifteen rounds exist. This argument is proved from
the executable rules in \texttt{ResidualMass.lean}.

A new fifteen-round cascade uses only 2-spawns after the original
$32{,}766$-round primed prefix. It attains both the mass bound and the
move bound. The bridge construction transports this cleaner suffix to
every opening. For an opening of mass at most 6, the same threshold
inequality used earlier gives $131072\le6+4(T-15)$ and hence
$T\ge32{,}782$; mass 8 gives $32{,}781$.
\end{proof}

The new minimum-mass witness uses $32{,}768$ spawned 4s and fifteen
spawned 2s, opening included, and scores $1{,}966{,}148$. Its score is
independently replayed, not asserted by the Lean mass theorem.

\subsection{Maximum score among the fastest games}

\begin{theorem}[Fastest-game score optimum]\label{thm:fastscore}
Among ordinary $4\times4$ plays reaching $131072$ in exactly
$32{,}781$ rounds, the largest score at that time is $1{,}966{,}216$.
\end{theorem}
\begin{proof}
The minimum-time equality argument forces all $131072$ mass at round
$32{,}766$ to contribute to the target. All earlier births are 4s, and
none of the final fifteen spawns can contribute to its ancestry.
The target's merge tree therefore contributes
\[
 (17-1)131072-4\cdot32768=1{,}966{,}080.
\]
Only the other fourteen births before the last spawn can contribute
additional merge score; the last spawn itself is unscored. If $y$ of
these fourteen births are 4s, their total mass is $28+2y$. Repeatedly
merging equal values in a geometry-free multiset increases its potential
and ends at the binary expansion of its mass. Write $\Phi_{\rm bin}(s)$
for that expansion's potential. The additional score is at most
\[
 \Phi_{\rm bin}(28+2y)-4y,\qquad 0\le y\le14.
\]
For these fifteen values of $y$, the bounds are respectively
\[
\begin{gathered}
68,64,120,116,116,112,120,116,\\
116,112,136,132,132,128,136.
\end{gathered}
\]
The maximum is 136. A second new fifteen-round cascade uses all 4s
and ends with the target together with junk tiles $32,16,8,4$,
attaining $1{,}966{,}080+136$. Both independent rule replayers accept
the complete trajectory in \texttt{witness/131072-fold-4.txt}.
\end{proof}

This is an optimum subject to the shortest-game constraint, not the
unrestricted $3{,}932{,}100$ score conjecture. Replacing its last spawn
by a 2 preserves the score. In particular the original minimum-length
certificate was not optimized for either endpoint mass or score:
\begin{center}\small
\begin{tabular}{@{}lrrr@{}}
\toprule
trajectory & rounds & endpoint mass & score\\
\midrule
original minimum-length witness & 32{,}781 & 131{,}116 & 1{,}966{,}092\\
new minimum-mass continuation & 32{,}781 & 131{,}102 & 1{,}966{,}148\\
new fastest-game score optimum & 32{,}781 & 131{,}132 & 1{,}966{,}216\\
\bottomrule
\end{tabular}
\end{center}

\paragraph{A second Pareto endpoint: one spawned 4.}
The gateway theorem below forces at least one 4-birth. Appending a new
all-2 cascade to the first gateway of the score-oriented witness gives
$131072$ in exactly $65{,}548$ rounds with just one 4-birth, opening
included, mass $131102$, and score $2{,}097{,}216$. This is the fewest
rounds possible under that resource constraint, because
$131072\le4+2(T-15)+2$ forces $T\ge65{,}548$.
It also has the largest score among minimum-mass target endpoints:
the non-target mass 30 has binary potential at most 68, and the required
4-birth costs at least four points, giving
$\Phi(131072)+68-4=2{,}097{,}216$.
Both existing replayers check \texttt{witness/131072-one-four.txt}.
This low-4 prefix is independently replayed, not included in the extended
Lean certificate.

\subsection{Unavoidable mass gateways}

\begin{theorem}[Gateway theorem]\label{thm:gateways}
Let $c\ge3$, $C=2^{c+2}$, and $G_j=C-2^{j+2}$ for
$j=0,\ldots,c-1$. Every ordinary play reaching mass at least $G_j$
visits mass exactly $G_j$, entering it by a 4-spawn. Its tile multiset
there is
\[
 \bigl(\{4,8,\ldots,2^{c+1}\}\setminus\{2^{j+2}\}\bigr)\cup\{4\},
\]
where the final union adds one tile, including a second 4 when $j>0$.
\end{theorem}
\begin{proof}
The binary expansion of $G_j-2$ has $c$ nonzero bits. Any representation
by at most $c$ powers of two is consequently that expansion itself:
the board is full, its values are distinct, and it is dead. A first
crossing of $G_j$ cannot depart from $G_j-2$, nor can it jump over
$G_j$. It must go from $G_j-4$ to $G_j$ by spawning a 4.
Immediately before that spawn there are at most $c-1$ tiles, because
one cell must be vacant, and $G_j-4$ has binary popcount $c-1$.
Its afterstate multiset is therefore the binary expansion, which fixes
the claimed post-spawn multiset. For $c\ge3$ every ordinary opening
lies below the first gateway.
\end{proof}

For $4\times4$ the gateways are, in order,
\[
\begin{gathered}
131072,196608,229376,245760,253952,258048,\\
260096,261120,261632,261888,262016,262080,\\
262112,262128,262136,262140.
\end{gathered}
\]
Thus \emph{every} route to $131072$ passes through the primed multiset,
not merely the fastest routes. Every full-chain play needs a 4-spawn
at each of the sixteen gateway entries. The argument is independent of
any snake invariant or merge-tree scheduling theorem.

\subsection{Stochastic consequences of the gateways}

Let each birth independently be 4 with probability $p$ and 2 with
probability $q=1-p$. In units of 2 the birth-mass process is a walk
with increments 1 and 2. Its probability of ever hitting $n\ge0$ is
\[
 u_n=\frac{1-(-p)^{n+1}}{1+p},
\]
as follows from $u_0=1$, $u_1=q$, and
$u_n=q u_{n-1}+p u_{n-2}$. A gateway at even distance $D$ can be
crossed only by hitting $D-2$ and next adding 2, not by hitting the
preceding dead layer. Its mass-only success probability is
\[
 b(D)=p u_{D-2}=\frac{p}{1+p}(1+p^{D-1}).
\]
Extending the independent birth sequence beyond an actual game's
stopping time gives an unrestricted walk containing every successful
game as an event. Discarding geometry and other failure conditions can
only increase the success probability. This coupling proves the
following bounds for every history-dependent player policy.

\begin{corollary}[Policy-independent probability bounds]\label{cor:probabilitygates}
On a $c$-cell board with $c\ge3$,
\[
 \Pr(\text{reach }2^{c+1})
 \le\frac{p}{1+p}(1+p^{2^c-1}),
\]
and
\[
 \Pr(\text{reach the full chain})
 \le\left(\frac{p}{1+p}\right)^c
       \prod_{k=1}^c(1+p^{2^k-1}).
\]
\end{corollary}
\begin{proof}
The first event requires crossing the first gateway, whose unit-mass
distance from zero is $2^c$. Successive distances for the full chain
are $2^c,2^{c-1},\ldots,2$; future increments are independent of the
history upon reaching each gateway. Multiply the corresponding $b(D)$.
The two initial tiles are the first two births, both before the first
barrier for $c\ge3$.
\end{proof}

For standard $4\times4$ play, with $p=1/10$, these become
\[
 \Pr(131072)\le\frac{1+10^{-65535}}{11}<0.090909091
\]
and a full-chain bound of approximately $2.3963146538\cdot10^{-17}$.
The exact expression in the corollary is the upper bound; a rounded
decimal must not silently replace it. In particular, the positive
correction in the first bound cannot be dropped to claim $\le1/11$.

There is also a better constructive lower bound than the original
minimum-length transcript provides. The one-4 witness first
reaches $131072$ at round $65{,}548$. Its exact rational probability,
including the prescribed two-cell opening and every empty-cell spawn
choice, is at least $10^{-48607}$. The computation factors that rational
number and proves the comparison by integer arithmetic. Hence
\[
 10^{-48607}\le \sup_\pi\Pr_\pi(131072)
 \le\frac{1+10^{-65535}}{11}.
\]
This interval is wide and is not an estimate of an optimal agent's
win rate. The new lower bound improves the valid $10^{-57048}$ bound
from the original shortest witness by 8441 powers of ten. The earlier
printed $10^{-57047}$ rounded that witness's probability in the wrong
direction; its logarithm is approximately $-57047.3719$.

\subsection{The value word of a ceiling-score game}

The score equality case has two post-spawn alternatives: the full chain
with $c$ spawned 4s, or the same chain with its last 4 replaced by a 2
and $c-1$ spawned 4s. The last slot in the score inequality contributes
zero in either case. Every earlier slot must be saturated. Changing
the final spawn from 4 to 2 changes no merge score.

The gateway theorem forces all $c-1$ nonfinal 4-births in either
ceiling-score play. There can be no other pre-final 4-births, so all
remaining births, including the opening, are 2s. Therefore the spawn
\emph{values} are fixed up to an arbitrary final 2 or 4: the unresolved
search need choose only directions and empty cells, not additional
2-versus-4 decisions. There are $2^{c+1}-2c-2$ forced 2-births and
$c-1$ forced 4-births before the final arbitrary birth. Consequently,
\[
 \Pr(\text{attain }\Phi_c-4c)
 \le q^{2^{c+1}-2c-2}p^{c-1}.
\]
For $4\times4$ this is
$0.9^{131038}0.1^{15}\approx1.0714006617\cdot10^{-6011}$.
Rip\`a already gave this value-word probability and the numerical
score and length formulas in 2014 \cite{ripa2014,ripaoeis2014}; we do
not claim them as new. This necessary-event bound does not assume that the score ceiling is
attainable. Our bounded fixed-value-word search did not complete a
ceiling-score witness and supplies no impossibility result.

\subsection{New full-chain arrangements}

An exhaustive 64-round suffix search from round $65{,}469$ of the
existing all-4 full-chain witness gives 32 distinct terminal $D_4$
orbits. Each complete chain has sixteen distinct values, hence trivial
stabilizer and eight labeled boards per orbit. A separate 70-round
checkpoint bridge supplies eight further noncorner-largest-tile orbits.
Deduplicating the union gives 40 orbits and at least
\textbf{320} arrangements with $65{,}533$-round witnesses, all optimal
by Corollary~\ref{cor:chain}. It does not classify all reachable endpoints:
only continuations of that fixed prefix were enumerated.

In particular, the question whether the 4 can be diagonally opposite
$131072$ has a positive answer. A sixteen-round replacement tail gives
\[
\begin{array}{rrrr}
4&8&16&32\\
64&128&256&512\\
8192&4096&2048&1024\\
16384&32768&65536&131072
\end{array}.
\]
The full transcript is \texttt{witness/full-chain-opposite.txt}.
The off-corner question also has a positive answer. A 70-round legal
bridge between rounds $49{,}166$ and $49{,}236$, followed by a reflected
suffix with the two distinct inert largest labels exchanged, produces
\[
\begin{array}{rrrr}
4&8&16&32\\
512&256&128&64\\
1024&2048&4096&8192\\
65536&131072&32768&16384
\end{array}.
\]
The explicit transcript \path{witness/full-chain-noncorner.txt} is
replayed in its entirety; no arbitrary label exchange is permitted as
an operation in the actual game. The exchange describes how the
candidate suffix was found, not how its legality is justified.
Symmetries put the largest tile in every boundary cell. Interior
placement in a full-chain endpoint and a complete classification remain
unresolved. Forty-eight representative histories (including overlaps)
are independently replayed by both rule implementations, and the
endpoint union is explicitly deduplicated. Reports are retained under
\path{research/results/}.

\subsection{A uniform one-dimensional family and a valid embedding}

On a line with $c\ge2$ cells, start with 4s at opposite ends and
repeatedly slide left and spawn a 4 at the far right. Nonzero values
remain nonincreasing. A non-full board moves because the last birth is
at the right edge; a full board with a duplicate has an adjacent merge.
The only full, distinct multiset with values at least 4 under the tile
ceiling is the full chain. Thus it is reached in $2^c-3$ rounds.
The first primed configuration folds by one carry per round, giving the
largest tile in $2^{c-1}+c-3$ rounds, exactly the general lower bound.

For maximum score, start with opposite-end 2s and spawn a 2 except
when the after-slide mass is $G_j-4$, when a 4 is required. Those
exceptional afterstates contain no 2s, so order remains nonincreasing.
A full distinct board containing a 2 must be at a dead layer $G_j-2$;
the rule skips every such layer. It ends on the full chain with exactly
$c$ 4-births, attaining both $\Phi_c-4c$ points and
$2^{c+1}-4-c$ rounds. This is an infinite-family construction, not merely
the finite tests (lengths 2 through 16) supplied with the paper.
It is not a construction for arbitrary rectangles.

\begin{proposition}[Floating-window embedding]\label{prop:embedding}
A legal play on a box of side lengths $n_i$ embeds in any box of side
lengths $N_i\ge n_i$, preserving its directions, spawn values, number
of rounds, and merge score, with translated spawn coordinates.
\end{proposition}
\begin{proof}
Keep all tiles in a translated $n_1\times\cdots\times n_d$ window.
After sliding along axis $j$, set that window's offset on axis $j$ to
0 toward the low wall or $N_j-n_j$ toward the high wall. Other offsets
are unchanged. Empty padding changes no ordered nonzero line sequence
and hence no merges; compression gives the translated small-board
result. The window moves in the same direction as the slide and cannot
cancel a nontrivial displacement. A merge likewise stays nontrivial.
Translate the prescribed spawn into its still-empty window cell.
Induction over the play proves the claim.
\end{proof}

Thus a fixed-coordinate sub-board is not isolated, but a moving-window
invariant does justify embedding. This repairs that part of the concern
raised about Das and Paul's induction; it does not prove their further
claim that every new cell buys an additional doubling. Independent lifts
of the entire shortest witness to $5\times5$, $5\times7$, and $4\times8$
check the invariant move by move.

\subsection{A computable upper bound on optimal expected score}
\label{sec:massbellman}

Use mass units of 2, put $H=2^{c+1}-2$, and let $F(m)$ be the
potential of the binary expansion of actual mass $2m$. For $c\ge3$,
there is a Bellman relaxation with only $H+1$ scalar mass states.
Stop the increment walk at its first mass with $c$ binary ones. Define
\[
W(m)=
\begin{cases}
F(m),&s_2(m)=c,\\
qW(m+1)+p\bigl(W(m+2)-4\bigr),&s_2(m)<c,
\end{cases}
\qquad 0\le m\le H.
\]
All successors used in the second case are at most $H$. This backward
recurrence exactly computes the expected terminal binary potential minus
four times the number of future 4-spawns in the relaxation.

To justify it as an upper bound, at spawn index $n$ let $m_n$ be the
mass and $f_n$ the number of 4-spawns, and put $R_n=F(m_n)-4f_n$.
The actual score is at most $R_n$. Moreover $R_n$ is nondecreasing:
a 2-spawn changes it by $F(m+1)-F(m)\ge0$, and a 4-spawn changes it by
$F(m+2)-F(m)-4\ge0$. The real game cannot continue beyond the walk's
first terminal barrier. Extending the pre-sampled walk after an earlier
geometric loss therefore still gives an upper bound. In particular,
for an ordinary random opening,
\[
\sup_\pi\mathbb E_\pi[\mathrm{score}]\le W(0).
\]
From an already given board $B$ of mass $m$ in units of 2, the analogous
upper bound on future score is $W(m)-\Phichain(B)$.
A second recurrence $J(m)=1+qJ(m+1)+pJ(m+2)$, with zero boundary values,
bounds the expected game length by $J(0)-2$.

For $c=16$, outward-rounded decimal evaluation gives
\begin{align*}
\sup_\pi\mathbb E_\pi[\mathrm{score}]&<1{,}915{,}854.958,\\
\sup_\pi\mathbb E_\pi[T]&<62{,}412.330.
\end{align*}
For comparison the four-cell score relaxation is $69.814912\ldots$,
above the exact geometric $2\times2$ value in Table~\ref{tab:honest}.
These are bounds, not strong solutions. The script
\path{research/mass_interval.py} uses fifty-digit arithmetic rounded
outwards at every operation and records enclosing intervals; a separate
forward computation checks the values numerically.

\subsection{Near-optimum rigidity and deadline probabilities}
\label{sec:slack}

If a $4\times4$ play reaches 131072 at round $T=32781+d$, inspect its
prefix ending fifteen moves earlier. The threshold-mass lemma requires
at least 131072 mass there. If $z$ is the number of 2-spawns in that
prefix, excluding the opening, and $M_0$ is the opening mass, the ledger
therefore gives
\[
2z+(8-M_0)\le4d.
\]
The left side also counts the opening 2s, at a charge of two mass units
each. Thus among the first $32768+d$ spawns, including the opening,
there can be at most $2d$ twos. The deterministic prefix inequality,
its general-target version, and the zero-slack consequence are
kernel-checked in \texttt{Game2048/DeadlineSlack.lean}.

Pre-sampling values again gives a bound for \emph{any} policy reaching
the tile by the deadline:
\[
\Pr_\pi(T_{131072}\le32781+d)
\le \sum_{z=0}^{\min(2d,32768+d)}
 \binom{32768+d}{z}(0.9)^z(0.1)^{32768+d-z}.
\]
For an earlier hit the required mass has already been acquired by the
specified prefix, so the same necessary event applies. At $d=0$ this
is $10^{-32768}$. For deadlines $40000$, $50000$, and $55000$, the
base-ten logarithms of the upper bounds are approximately
$-14854.03$, $-3668.61$, and $-909.55$. These are spawn-value bounds;
requiring suitable locations can only decrease the success probability.
The stochastic corollary is not separately formalised in Lean.

\paragraph{Reproduction and scope.}
The scripts in \texttt{research/} generate the short bridges, clean
cascades, suffix variants, exact probability bounds, and line tests.
\texttt{bash research/finish.sh} runs the extended checks after the
original formal development is rebuilt. The symmetry, every-opening,
endpoint-mass, and deterministic deadline-slack theorems are kernel-checked. The other new deductions
have written proofs and executable checks, and are not thereby claimed
to have been completely formalized in Lean.
